\subsection{An explicit two-unit family with either local behavior at the residue prime}
\label{sec:kt-explicit-two-unit}

The following is a partial solution of the simultaneous Bernoulli-unit
question.  It gives two of its three units, at a fixed prime and in an
unbounded weight progression, together with either splitting behavior at
that prime.  The third unit and any additional local conditions are not
asserted.

\begin{theorem}[KT-EXPLICIT-TWO-UNIT-FAMILY]
Let $\ell\geq7$ be a prime with $\ell\equiv3\pmod4$.  For each sign
$s\in\{1,-1\}$, write
\[
 \ell+2s=m_su_s^2,\qquad m_s>0\text{ square-free},\qquad
 D_s=-4m_s,\qquad E_s=\mathbf Q(\sqrt{D_s}).                 \tag{EU.1}
\]
Then $D_s$ is a negative fundamental discriminant and $\ell\nmid D_s$.
For every integer $t\geq0$, put
\[
 e=\frac{\ell-3}{2}+(\ell-1)t.                              \tag{EU.2}
\]
This is a positive even integer satisfying $(\ell-1)\nmid2e$.
The discriminants here are even; as explained below, this arithmetic
family does not satisfy the source's full GD condition when $N$ is even.
With the fixed $\ell$-adic Teichm\"uller character $\omega$, both
\[
 B_{1,\chi_{D_s}}\quad\text{and}\quad
 B_{1,\chi_{D_s}\omega^e}
                                                               \tag{EU.3}
\]
are $\ell$-adic units.  The normalized Bernoulli number
$B_{e+1,\chi_{D_s}}/(e+1)$ is therefore also a unit.
Moreover,
\[
 \chi_{D_+}(\ell)=-\left(\frac2\ell\right),\qquad
 \chi_{D_-}(\ell)= \left(\frac2\ell\right).                 \tag{EU.4}
\]
Thus one of the two fields is split at $\ell$ and the other is inert.
The two nontrivial Teichm\"uller twists $\omega^e$ and $\omega^{2e}$
have order $(\ell-1)/2>2$.

For either sign, the remaining third-unit condition is exactly
\[
 B_{1,\chi_{D_s}\omega^{2e}}\in\mathbf Z_\ell^\times
 \quad\Longleftrightarrow\quad
 \frac{B_{\ell-2,\chi_{D_s}}}{\ell-2}\in\mathbf Z_\ell^\times.
                                                               \tag{EU.5}
\]
Neither side of (EU.5) is claimed by this theorem.
\end{theorem}

\begin{proof}
Both $\ell+2$ and $\ell-2$ are positive odd integers congruent to $1$
modulo $4$.  Their square-free parts $m_s$ therefore also satisfy
$m_s\equiv1\pmod4$.  The integer $-m_s$ is square-free and congruent
to $3$ modulo $4$, so its quadratic field has fundamental discriminant
$D_s=-4m_s$.  This includes the square cases $m_s=1$, when $D_s=-4$.
Since $\ell$ divides neither $\ell+2$ nor $\ell-2$, it does not divide
$D_s$.  Reduction of (EU.1) modulo $\ell$ gives (EU.4), because
$(-1/\ell)=-1$ and $u_s^2$ is a nonzero square modulo $\ell$.

We first check the ordinary class-number factor, including all small
cases.  Theorem~4.1 on PDF p.~14 of Benmerieme--Movahhedi gives,
since $\chi_{D_s}(2)=0$,
\[
 h(D_s)=\frac{w(E_s)}4\sum_{r=1}^{2m_s}\chi_{D_s}(r).
                                                               \tag{EU.6}
\]
Only the $m_s$ odd terms can be nonzero, and each is at most $1$.
If $m_s>1$, then $w(E_s)=2$ and hence
$h(D_s)\leq m_s/2\leq(\ell+2)/2<\ell$.
If $m_s=1$, the same formula gives $h(-4)=1<\ell$.
Therefore $\ell\nmid h(D_s)$, including both signs at $\ell=7$.
The imaginary quadratic class-number formula
$B_{1,\chi_{D_s}}=-2h(D_s)/w(E_s)$, where $w(E_s)=2$ or $4$,
proves the first assertion in (EU.3).

Here is the imported arithmetic input for the second assertion.
Benmerieme--Movahhedi prove that
\[
             F_s=\mathbf Q(\sqrt{\ell(\ell+2s)})
                                                               \tag{EU.7}
\]
is $\ell$-rational for every prime $\ell>3$: their Proposition~3.2
gives the plus sign, and the proof of Proposition~4.4 gives the minus
sign.  Their Corollary~2.2 identifies this with
$L_\ell(1,\chi_{F_s})\in\mathbf Z_\ell^\times$.
Equation~(3.2) then gives, for $n_0=(\ell-1)/2$,
\[
 L_\ell(1,\chi_{F_s})
       \equiv L(1-n_0,\chi_{D_s})
       =-\frac{B_{n_0,\chi_{D_s}}}{n_0}\pmod\ell.          \tag{EU.8}
\]
Indeed the quadratic character of $\mathbf Q(\sqrt{-\ell})$
is $\omega^{n_0}$, and the primitive character underlying
$\chi_{F_s}\omega^{n_0}$ is $\chi_{D_s}$.  This is also equality
as Galois characters.  As an imprimitive Dirichlet function the
product on the left still omits the $\ell$ Euler factor; the
interpolation value therefore includes
$1-\chi_{D_s}(\ell)\ell^{n_0-1}$, which is $1$ modulo $\ell$
because $n_0\geq3$.  Thus it does not change (EU.8).
The hypotheses of that
displayed congruence hold for the square-free negative integer $-m_s$:
it has the prescribed sign, is prime to $\ell$, and is different from
$-\ell$.  Thus $B_{n_0,\chi_{D_s}}/n_0$ is a unit.

Now $n_0$ is odd, $3\leq n_0\leq\ell-2$, and
$e+1\equiv n_0\pmod{\ell-1}$.  Applying the already proved
KT-FIELD-CHARACTER-REDUCTION yields
\[
 \frac{B_{e+1,\chi_{D_s}}}{e+1}
 \equiv\frac{B_{n_0,\chi_{D_s}}}{n_0}
 \equiv B_{1,\chi_{D_s}\omega^e}\pmod\ell.                 \tag{EU.9}
\]
Its integrality statement applies even when the original weight
$e+1$ is divisible by $\ell$ or exceeds $\ell$.  It also proves
(EU.5), since $2e+1\equiv\ell-2\pmod{\ell-1}$.

Finally, $2e\equiv-2\pmod{\ell-1}$ proves the required GD
nondivisibility.  Writing $\ell=4a+3$ gives
$\gcd(\ell-1,e)=\gcd(4a+2,2a)=2$, so $\omega^e$ has order
$(\ell-1)/2$, an odd integer at least $3$.  Squaring does not change
this order, which proves the final assertion concerning the twists.
\end{proof}

\paragraph{Exact primary locators and check scope.}
Y.~Benmerieme and A.~Movahhedi, \emph{Multi-quadratic $p$-rational
number fields}, Journal of Pure and Applied Algebra 225 (2021),
106657, DOI \texttt{10.1016/j.jpaa.2020.106657}.
The version used here is the authors' arXiv:2007.04864v1,
9 July 2020, \url{https://arxiv.org/pdf/2007.04864v1}:
Corollary~2.2 on PDF p.~5 and its preceding derivation on pp.~4--5;
Proposition~3.2 and its entire proof on pp.~11--12;
equation~(3.2) and its explanatory paragraph on p.~12;
Theorem~4.1 on p.~14; Proposition~4.4 and its entire proof on p.~15.
The present proof imports those arithmetic statements, not fresh
proofs of their cited Coates or $p$-adic $L$-function results.
The first class-number unit follows directly from the finite character
sum (EU.6), with both square cases included.
This theorem does not prescribe local behavior at any prime other than
$\ell$, does not prove (EU.5), and does not exclude every stable
automorphic congruence even if that remaining unit is later verified.

\begin{proposition}[KT-EXPLICIT-TWO-UNIT-GD-SCOPE]
Neither discriminant in (EU.1) is an admissible auxiliary discriminant
for the actual source GD condition with even $N$, in particular with
$N=2$.
\end{proposition}
\begin{proof}
The actual source hypotheses \texttt{GD.N\_even} and
\texttt{GD.cop} state that $N$ is even and coprime to the absolute
discriminant of the auxiliary quadratic field.  Thus $2\mid N$
and $\gcd(N,|D|)=1$ imply $2\nmid D$.  This exact deduction is
already written in the pinned source
\texttt{UnitaryThreeGroups/AdmissibleOfTypeII.lean:119} as
\texttt{Nat.coprime\_two\_right.1
(D.gd.cop.coprime\_dvd\_right D.gd.N\_even.two\_dvd)}.
But (EU.1) has $D_s=-4m_s$, divisible by four for both signs,
including $m_s=1$.  Therefore the family proves the two asserted
Bernoulli units and the exponent nondivisibility, but supplies no
field satisfying the full GD hypotheses.  The bounded third-unit
certificate for this family has the same parity limitation.
\end{proof}
